% =============================================================================
\section{\sys: Learning From What the Kernel Already Sees}
\label{sec:design}
% Budget: 3.0 pages. Draft SECOND (Step 5.2). One subsection per challenge, each
% "Observation. -> Solution." Misprediction handling is its own subsection.
% Self-check: a reader can follow ONE eviction step by step; the worst case under a
% bad model is stated; every named component appears in the evaluation.
% =============================================================================

\subsection{Overview and Formulation}
\label{sec:design-overview}

\plan{Walk Figure~\ref{fig:overview} in decision order.
\emph{Offline}: complete reference strings $\rightarrow$ replay under LRU at 5/10/20\%
memory $\rightarrow$ record eviction decisions (candidates, features, label)
$\rightarrow$ fit a linear scorer $\rightarrow$ select feature subset and probation on
validation only $\rightarrow$ export integer weights.
\emph{Online}: at each eviction the kernel scans the eligible resident pages, updates
their counters, scores each with an integer dot product, protects pages in probation,
and evicts the highest score (oldest load on a tie).}

\plan{Formulation, and why this kind of ML. The policy imitates \belady: for each
candidate it regresses $\log(1+d)$, where $d$ is the distance in references to the
page's next use (capped at $2^{20}$), and evicts the candidate with the largest
prediction. Supervised imitation fits because complete reference strings give exact
labels offline and \belady is optimal for faults; reinforcement learning would need
exploration inside a running kernel. Decisions compound, so every policy is judged by
replaying complete streams, never by per-reference accuracy; on-policy data (DAgger)
was tried and did not help (\S\ref{sec:eval-ablate}). Why linear: five heavier models,
trained the same way and also run in the kernel, were level with it at 20--150$\times$ the
cost (\S\ref{sec:eval-ablate}). The global model is trained on six workloads (B-tree,
graph, key-value, LZW, matmul, sort) and used unchanged on all others.}

\begin{figure*}[t]
\centering
\figboxwide{1.7in}{Overview. Left (offline, host): traced workload $\rightarrow$
reference string with R/W marks $\rightarrow$ LRU replay $\rightarrow$
decision-point dataset (\belady label, recency-stratified candidates) $\rightarrow$
ridge regression $\rightarrow$ nested selection on validation (subset $\times$
probation) $\rightarrow$ integer weights (\texttt{mlweights.h} / \texttt{vmctl}).
Right (online, kernel): fault $\rightarrow$ scan candidates (read/clear A, update
counters) $\rightarrow$ features in Q8 $\rightarrow$ $\sum_j q_j X_j$ $\rightarrow$
probation $\rightarrow$ victim (optionally among 32 sampled candidates).}
{to draw; single vector figure, the paper's mechanism figure}
\caption{\plan{Takeaway: training is offline and uses the future; the kernel side is a
scan it already performs plus a few integer multiply-adds per page.}}
\label{fig:overview}
\end{figure*}

% -----------------------------------------------------------------------------
\subsection{Features the Kernel Already Keeps}
\label{sec:design-features}
% Challenge C1 -> technique T1 -> ablation A1 (tier comparison, Table tab:tiers)

\finding{Observation.} A kernel never sees individual accesses
(\S\ref{sec:bg-mechanism}), so a model trained on per-access recency or frequency cannot
be deployed as trained.

\finding{Solution.} \sys uses only features a kernel can compute when it scans the
resident pages at an eviction (tier \tierK), plus refault bookkeeping that Linux already
keeps in its workingset shadow entries (tier \tierKp) (Table~\ref{tab:features}). Time is
counted in scans, one per eviction. Counts enter as $\log(1+x)$ and are standardised.
Oracle features (tier \tierF: exact recency, frequency, stack distance, write ratio)
serve only as an ablation that bounds what extra information is worth.

\begin{table}[t]
\caption{Kernel-observable features. All are small per-page integers; the last column
names the Linux state that already holds an equivalent.}
\label{tab:features}
\vskip 0.1in
\begin{center}
\begin{footnotesize}
\setlength{\tabcolsep}{3pt}
\begin{tabular}{@{}l>{\raggedright\arraybackslash}p{1.2in}>{\raggedright\arraybackslash}p{1.15in}@{}}
\toprule
Feature & Definition & Linux analogue \\
\midrule
\multicolumn{3}{@{}l}{\textit{\tierK: from accessed/dirty-bit scans}}\\
\feat{ref}      & A bit at this scan                     & young bit, PG\_referenced \\
\feat{aging}    & 8-bit aging counter                    & -- \\
\feat{sfreq}    & scans that found it accessed           & DAMON access count; MGLRU tiers \\
\feat{idle}     & scans since last found accessed        & MGLRU generation age; DAMON age \\
\feat{age}      & scans since it was loaded              & -- \\
\feat{dirty}    & D bit                                  & PG\_dirty \\
\midrule
\multicolumn{3}{@{}l}{\textit{\tierKp: refault bookkeeping}}\\
\feat{refaults} & times evicted and faulted back         & workingset shadow entry \\
\feat{rdist}    & evictions between its eviction and refault & workingset refault distance \\
\bottomrule
\end{tabular}
\end{footnotesize}
\end{center}
\vskip -0.1in
\end{table}
% src: ML_REPORT section 4 and 9. \verify{Linux version numbers for MGLRU (merged 6.1?) and DAMON (5.15?)}

% -----------------------------------------------------------------------------
\subsection{Training on Decisions, Not References}
\label{sec:design-train}
% Challenge C3 -> technique T3 -> ablations A3 (v1 vs v2; probation sweep; DAgger)

\finding{Observation.} A policy acts at an eviction, over the pages resident at that
moment, and a naive sample of those pages almost never contains the page just used
(\S\ref{sec:bg-naive}).

\finding{Solution.} \sys trains on recorded eviction decisions. Every training and
validation stream is replayed under LRU at 5, 10 and 20\% of its distinct pages; at a
sample of evictions we record up to 32 resident candidates with their features and
labels: always \belady's victim, always the four most recently accessed candidates
(recency-stratified recording), and the rest uniformly. This raises the share of
just-accessed rows from 0.6\% to 3.1\%. Second, a \emph{probation} rule: a page loaded
fewer than $p$ scans ago is not evicted while an older candidate exists, which is
\clock's second chance applied to a learned score. $p\in\{0,1,2,4\}$ is chosen per
scope on validation. \todo{training rows and decisions for the current six-workload
datasets (4.73M rows over 357K decisions before LZW was added and the sampling seed was
fixed)}
% src: ML_REPORT 5.2, 7.2; week-15 addendum (CRC32 seed; lzw in the pipeline)

% -----------------------------------------------------------------------------
\subsection{Integer-Only Scoring}
\label{sec:design-int}
% Challenge C2 -> technique T2 -> ablations A2 (float vs 4/8/12-bit; linear vs MLP)

\finding{Observation.} xv6 saves no floating-point registers on a context switch, and a
production kernel avoids floating point on the fault path.
\verify{cite the Linux policy on kernel FPU use and its cost}

\finding{Solution.} The score of a candidate page is
\begin{equation}
s(\text{page}) = \sum_{j\in S} q_j\, X_j(\text{page}), \qquad
q_j = \operatorname{round}\!\left(\frac{w_j}{\sigma_j}\, 2^{b}\right),
\label{eq:score}
\end{equation}
where $X_j$ is feature $j$ in Q8 fixed point, $S$ is the selected subset (1--5
features), and the standardisation is folded into the weight. Means and the bias are the
same for every candidate and drop out of the arg-max. $\log(1+x)$ comes from a 256-entry
table plus a most-significant-bit search and a 1,024-entry mantissa table, within 1.04
Q8 units of exact. The same header computes features in the kernel and in the
simulator, so a model scores a page identically in both. The kernel's default weights
are the global model over five features (\feat{idle}, \feat{age}, \feat{dirty},
\feat{refaults}, \feat{rdist}; probation 4 scans).

\finding{Heavier models, in whole numbers too.} To test whether capacity pays, five
heavier models (a CNN over the aging bits, a deeper MLP, set-attention, boosted trees and
ranking trees) are converted to whole-number arithmetic with a shared evaluator that the
kernel and the simulator both compile, and selected per process like the linear weights.
% src: ML_REPORT 8, 8b; kernel/mlweights.h (current default); kernel/mlnet.h; AUSTRA_HANDOFF week-15 addendum

% -----------------------------------------------------------------------------
\subsection{When the Model Is Wrong}
\label{sec:design-wrong}
% Challenge C4 -> technique T4 -> ablation A4 (deliberately wrong models). NON-NEGOTIABLE.
% Every accepted ML-for-OS paper has a design element here, not a sentence.

\plan{State the worst case first, then the layers that bound it.}
\begin{enumerate}
\item \textbf{Correctness never depends on the model.} The score only orders eligible
  candidates (the faulting process's own evictable frames). The kernel checks every
  returned victim against that set and falls back to \clock if it is not in it (counted
  as \texttt{policy\_fallbacks}). The same workload produces byte-identical output under
  every policy (data-invariance suite, \S\ref{sec:impl}). A bad model costs
  performance, not data. \todo{report policy\_fallbacks over the 2,375 runs (expected 0)}
\item \textbf{Probation} keeps the newest pages out of the choice, removing the failure
  of \S\ref{sec:bg-naive}.
\item \textbf{Validation-only selection.} Many feature subsets are catastrophic
  ($\geq3\times$ \clock on graph for 98--126 of 255 kernel subsets); feature subset and
  probation are chosen jointly on validation, never on test or held-out streams.
\item \textbf{A fallback guard (simulator prototype).} The learned policy and \clock
  alternate on the same memory: one window under each, compare their faults, use the
  better for the next $K$ windows, then test again. On streams the settings were not
  tuned on, it cuts the worst case from 1.86 to 1.36 of \clock's faults and the losing
  cells from 1.29 to 1.07, at a cost of about two points on average (0.861 to 0.881 of
  \clock). It is \emph{not in the kernel}: the window is measured in accesses, which a
  kernel cannot count; a kernel version needs another measure of progress.
\end{enumerate}
\plan{State plainly that the kernel runs have no guard, so the failures of
\S\ref{sec:eval-loss} are what a deployment would see today.}
\todo{decision for the authors: port the guard to the kernel (window in evictions or
faults), or keep it as a prototype and say so; the collaborator calls it an optional extra}
% src: week-15 report, "A safety guard (prototype, simulator only)"; AUSTRA_HANDOFF START HERE section 2

% -----------------------------------------------------------------------------
\subsection{Bounding the Decision Cost}
\label{sec:design-cost}
% Challenge C5 -> technique T5 -> ablation A5 (sample size k vs faults and instructions).

\finding{Observation.} Scoring every eligible page costs about 340 instructions per
candidate, so the decision grows in a straight line with the resident set
(\S\ref{sec:eval-cost}).

\finding{Solution.} Score only $k$ randomly sampled candidates, drawn with a partial
Fisher--Yates shuffle over the candidate array (no extra memory). With $k=32$ the
decision costs about 11,700 instructions whatever the memory size, for 0 to 1 point of
faults; $k=16$ costs 5,850 instructions for 1 to 2 points. The weights are unchanged;
$k$ is a per-process setting.
% src: week-15 report, "In the xv6 kernel" (lzw 0.837->0.842, Twitter replays 0.876->0.875,
% patbench 0.853->0.861 at k=32); week_16 report section "What choosing a victim costs"

% -----------------------------------------------------------------------------
\begin{algorithm}[t]
\caption{\sys victim selection (\texttt{choose\_ml}), one eviction}
\label{alg:choose}
\begin{algorithmic}
\STATE {\bfseries Input:} eligible resident pages $C$, weights $q$, subset $S$, probation $p$, scan counter $t$
\STATE {\bfseries if} sample size $k$ is set {\bfseries then} $C \leftarrow k$ random pages of $C$ \hfill (partial Fisher--Yates)
\STATE $old \leftarrow$ false; \ $best \leftarrow \bot$
\FOR{each page $c \in C$}
  \STATE read and clear $A(c)$; update \feat{aging}, \feat{sfreq}, last-seen scan
  \STATE {\bfseries if} $t - \mathrm{load}(c) \ge p$ {\bfseries then} $old \leftarrow$ true
\ENDFOR
\STATE flush the TLB once if any A bit was cleared
\FOR{each page $c \in C$}
  \IF{$old$ {\bfseries and} $t - \mathrm{load}(c) < p$}
    \STATE skip $c$ \hfill (probation)
  \ENDIF
  \STATE $s_c \leftarrow \sum_{j\in S} q_j X_j(c)$ \hfill (integers, Q8)
  \STATE keep $c$ if $s_c > s_{best}$, or equal and loaded earlier
\ENDFOR
\STATE {\bfseries return} $best$ \hfill (kernel validates it; \clock if invalid)
\end{algorithmic}
\end{algorithm}
% src: kernel/vmpage.c choose_ml() (39 lines) and choose_policy_victim(); kernel/mlfeat.h (63 lines)
